% Part: normal-modal-logic
% Chapter: axioms-systems
% Section: consistency

\documentclass[../../../include/open-logic-section]{subfiles}

\begin{document}

\olfileid{nml}{prf}{con}

\olsection{সঙ্গতি}

সঙ্গতি !!{formula}s-এর সমষ্টির একটি গুরুত্বপূর্ণ বৈশিষ্ট্য। নির্দিষ্ট পদ্ধতিতে কোনো !!{formula}s-এর সমষ্টি থেকে~$\lfalse$-এর মতো স্ববিরোধ !!{derivable} হলে সমষ্টিটি সেই পদ্ধতির সাপেক্ষে অসঙ্গত; নইলে সঙ্গত। সমষ্টিটি অসঙ্গত হলে তার সব !!{formula}s সেই পদ্ধতির কোনো মডেলের একটি জগতে একসঙ্গে সত্য হতে পারে না। এখানে পদ্ধতির মডেল বলতে এমন মডেল বোঝায়, যার প্রত্যেক জগতে পদ্ধতির সমস্ত স্বতঃসিদ্ধের প্রতিস্থাপনজাত উদাহরণ সত্য। পূর্ণতার উপপাদ্যের জন্য আমরা বিপরীতটিও প্রমাণ করব: সংশ্লিষ্ট পদ্ধতির সাপেক্ষে প্রতিটি সঙ্গত সমষ্টি সেই পদ্ধতির একটি মডেলের কোনো জগতে সত্য; সেই মডেলটি ‘প্রামাণিক মডেল’।
% BN-SRC-700 TARGET-CORRECTION-BEGIN
\emph{উৎস-সংশোধন:} অসঙ্গতি পদ্ধতি-সাপেক্ষ। উৎসের অনির্দিষ্ট ‘মডেল’ কথাটিতে সংশ্লিষ্ট পদ্ধতির স্বতঃসিদ্ধগুলি মেনে চলার শর্ত যোগ করা হয়েছে; অন্য পদ্ধতির মডেলে সূত্রগুলির যুগপৎ সত্যতা অসম্ভব বলা হয়নি।
% BN-SRC-700 TARGET-CORRECTION-END

\begin{defn}
  সমষ্টি $\Gamma$ পদ্ধতি~$\Sigma$-র সাপেক্ষে \emph{সঙ্গতি-সম্পন্ন}, অথবা সংক্ষেপে $\Sigma$-সঙ্গত, তখনই এবং কেবল তখনই, যখন $\Gamma \Proves/[\Sigma] \lfalse$।
\end{defn}

যেমন, সমষ্টি $\{ \Box(p \lif q), \Box p, \lnot\Box q \}$ প্রতিজ্ঞামূলক যুক্তিবিদ্যার সাপেক্ষে সঙ্গতি-সম্পন্ন, কিন্তু \Log{K}-সঙ্গত নয়। একইভাবে $\{ \Diamond p, \Box\Diamond p \lif q, \lnot q \}$ \Log{K5}-সঙ্গত নয়।

\begin{prop}\ollabel{prop:consistencyfacts}
  ধরা যাক $\Gamma$ !!{formula}s-এর একটি সমষ্টি। তবে:
  \begin{enumerate}
  \item $\Gamma$ তখনই এবং কেবল তখনই $\Sigma$-সঙ্গত, যখন এমন কোনো !!{formula}~$!A$ আছে যার জন্য $\Gamma \Proves/[\Sigma] !A$।
  \item \ollabel{prop:consistencyfacts-b}%
    $\Gamma \Proves[\Sigma] !A$ তখনই এবং কেবল তখনই, যখন $\Gamma \cup \{ \lnot!A \}$ $\Sigma$-সঙ্গত নয়।
  \item \ollabel{prop:consistencyfacts-c}%
    $\Gamma$ $\Sigma$-সঙ্গত হলে যেকোনো !!{formula} $!A$-এর জন্য হয় $\Gamma \cup \{ !A \}$ $\Sigma$-সঙ্গত, নয়তো $\Gamma \cup \{ \lnot!A \}$ $\Sigma$-সঙ্গত।
  \end{enumerate}
\end{prop}

\begin{proof}
  ধ্রুপদি প্রতিজ্ঞামূলক যুক্তিবিদ্যা থেকে এই তথ্যগুলি সহজেই পাওয়া যায়। \olref{prop:consistencyfacts-c}-এর যুক্তিটি দিই। বিপরীতক্রমে ধরি, $\Gamma \cup \{ !A \}$ এবং $\Gamma \cup \{ \lnot!A \}$—কোনোটিই $\Sigma$-সঙ্গত নয়। তখন অসঙ্গতির সংজ্ঞা অনুসারে $\Gamma, !A \Proves[\Sigma] \lfalse$ এবং $\Gamma, \lnot !A \Proves[\Sigma] \lfalse$—উভয়ই সত্য। নিষ্পাদন উপপাদ্য অনুসারে $\Gamma \Proves[\Sigma] !A \to \lfalse$ এবং $\Gamma \Proves[\Sigma] \lnot!A \lif \lfalse$। কিন্তু $(!A \lif \lfalse) \lif ((\lnot!A \lif \lfalse) \lif \lfalse)$ সর্বতঃসত্য সূত্রের একটি উদাহরণ। অতএব \olref[prp]{prop:derivabilityfacts}\olref[prp]{prop:derivabilityfacts-ruleT} থেকে $\Gamma \Proves[\Sigma] \lfalse$।
  % BN-SRC-701 TARGET-CORRECTION-BEGIN
  \emph{উৎস-সংশোধন:} অসঙ্গত সমষ্টি থেকে স্ববিরোধের প্রমাণযোগ্যতা এখানে সংজ্ঞা থেকেই আসে। উৎসে অপ্রয়োজনীয় ও সরাসরি এই ধাপটি না-বলা দ্বিতীয় অংশের নির্দেশ সরানো হয়েছে।
  % BN-SRC-701 TARGET-CORRECTION-END
\end{proof}

\end{document}
